\section{The Public Message and Its Two Strata}

\subsection{What a ``text'' of the discourse can be}

There is no autograph of the La Salette discourse. There is no
shorthand record, no contemporaneous transcript taken while the
children spoke, and no manuscript in either child's hand, because
neither could write. What exists is a recital that two illiterate
children repeated, first on the evening of 19 September 1846 and
then hundreds of times, partly in French and partly in the
\term{patois} of the Corps country, written down by adults who had
to translate as they wrote.

That is not a scandal; it is the ordinary condition of an oral
deposition. But it fixes the ceiling. The stable content of the
discourse is well attested by 1847 and can be inventoried with
confidence. Syllable-level inerrancy of any French sentence cannot
be, and this study never asserts it. Where a word differs between
strata---and one important word does---the difference is printed
rather than smoothed.

\subsection{The recital stratum: the discourse as an 1854 witness
prints it from the diocesan documents}

The most useful witness this project could reach for the early
stratum is the 1854 pilgrimage narrative of an English bishop who
travelled to Corps, read the diocesan documents, and printed the
children's recital in parallel columns---French where the Lady is
reported to have spoken French, \term{patois} where she is reported
to have repeated herself in \term{patois}, with an English rendering
beside each. The parallel setting is itself evidence: the witness
did not harmonize the two languages into one smooth French text, and
the seam where the Lady changes language survives in his page.

\begin{witnesstext}{Opening of the discourse, French column, as
printed in 1854}
«~Si mon peuple ne veut pas se soumettre, je suis forcée de laisser
aller la main de mon fils. Elle est si forte, si pesante, que je ne
peux plus la maintenir. Depuis le temps que je souffre pour vous
autres! Si je veux que mon fils ne vous abandonne pas, je suis
chargée de le prier sans cesse. Et pour vous autres vous n'en faites
pas cas.~»\endnote{Ullathorne 1854, pp.~36--37, French column of
the parallel text. The English column of the same page reads: ``If
my people will not submit, I shall be forced to let go the hand of
my Son. It is so strong, so heavy, that I can no longer withhold
it\ldots'' The 1854 book prints the French with the period's
lower-case ``fils''; the capitalization is reproduced as printed and
carries no theological weight.}
\end{witnesstext}

The \term{patois} seam occurs at the potatoes. Mélanie, who did not
know the French word \term{pommes de terre}, was about to ask
Maximin what it meant; the Lady said---in the 1854 witness's
English---``Ah! my children, you do not understand; I will say it in
a different way,'' and began again in the local speech: ``Si las
truffas se gastoun\ldots'' The recital then runs, still in
\term{patois}, through the corn that will be eaten by beasts and
fall to dust in the threshing, the great famine, the children under
seven who will be seized with trembling and die in the hands of
those who hold them, the walnuts that will go hollow and the grapes
that will rot, and the promise that if they are converted the stones
and rocks will become heaps of corn and the potatoes will be sown of
themselves.\endnote{Ullathorne 1854, pp.~38--39, \term{patois} and
English columns. \term{Las truffas} is the local word for potatoes,
the very word Mélanie knew and \term{pommes de terre} the one she
did not; the witness prints both. The prophecy of children dying is
reported here strictly as part of the recital and is not treated by
this study, or by the 1851 act, as a verified prediction or as an
event to be dated (\S3.5).}

Then the catechetical exchange: ``Do you say your prayers well, my
children?''---``Not very well, Madam''---and the instruction to say
them morning and evening, ``at least an Our Father and a Hail
Mary,'' and more when they can; then the complaint that only a few
old women go to Mass while the rest work Sundays all summer, that in
winter the lads go to Mass only to mock at religion, and that in
Lent ``they go to the shambles like dogs.'' Then the question about
spoilt corn, Maximin's denial, and the Lady's reminder of the day at
the hamlet of Le Coin when his father showed him ears of wheat that
crumbled to dust in their hands and gave him a piece of bread on the
road home, saying he did not know who would eat any next year if the
corn went on so---at which the boy answered, in \term{patois},
``Oh! yes, Madam, I recollect now; just this moment, I did not
remember.'' The discourse closes with the charge repeated twice in
French: ``Well, my children, you will make this to be known to all
my people.''\endnote{Ullathorne 1854, pp.~40--42. The Coin episode is
the discourse's one verifiable-in-principle element: it reports a
private incident between Maximin and his father that the child
himself had forgotten and, on the recital's own showing, recognized
only when reminded. This study reports it as part of the recital and
records that it did not reach any independent testimony of the
father.}

\subsection{The received custodial text}

The text now published as ``le Message de La Salette'' is a
harmonized French discourse without apparatus, with the
\term{patois} rendered into French and the moments when the Lady
spoke privately to each child marked by editorial stage directions
in parentheses. It is the text the Sanctuary and the Missionaries
put before pilgrims, and it is faithful in substance to the recital
stratum. Its opening and closing lines:

\begin{witnesstext}{Received custodial text, opening and closing
(Sanctuary of Notre-Dame de La Salette)}
«~Avancez, mes enfants, n'ayez pas peur, je suis ici pour vous
conter une grande nouvelle.\par
Si mon peuple ne veut pas se soumettre, je suis forcée de laisser
aller le bras de mon Fils. Il est si fort et si pesant que je ne
puis plus le maintenir.~[\ldots]\par
Eh bien, mes enfants, vous le ferez passer à tout mon peuple!
Allons, mes enfants, faites-le bien passer à tout mon
peuple!~»\endnote{Sanctuaire Notre-Dame de La Salette, ``Le Message
de La Salette,'' read 2026-07-25 in the Internet Archive capture of
5 November 2012 (References); the Missionaries of Our Lady of La
Salette publish the same received text
(\url{https://www.missionnaire-lasalette.fr/spip.php?article25}, read
2026-07-25). Only the opening and closing lines are reproduced here;
the whole received text is not reproduced in this study, and the
custodial publishers' pages are cited for it. The parenthetical
stage directions in the received text---``Jusqu'ici la Belle Dame a
parlé en français\ldots\ termine son discours en patois''; ``A ce
moment Mélanie voit que la Belle Dame dit quelques mots à Maximin,
mais elle n'entend pas''---are editorial in the received text
itself, not reported speech.}
\end{witnesstext}

\subsection{Element inventory: where the strata agree and where they
differ}

The comparison below runs element by element over the whole
discourse. It is the study's principal textual instrument, and its
result is unspectacular in the best way: the strata agree on
substance throughout, and differ in three places that a reader
should know about.

\begin{stratatable}
Opening address & ``Come near, my children, be not afraid, I am here
to tell you great news'' (English column) & «~Avancez, mes enfants,
n'ayez pas peur, je suis ici pour vous conter une grande nouvelle~»
& Agree \\
The Son's hand or arm & «~laisser aller \emph{la main} de mon fils.
Elle est si forte, si pesante~» & «~laisser aller \emph{le bras} de
mon Fils. Il est si fort et si pesant~» & \textbf{Differ}: hand
\slash\ arm, with the pronouns and adjectives following \\
Constant intercession & ``I am compelled to pray to Him without
ceasing''; ``you take no heed of it'' & «~je suis chargée de le
prier sans cesse~»; «~vous n'en faites pas cas~» & Agree \\
Unrepayable pains & ``However much you pray\ldots\ you will never
recompense the pains I have taken for you'' & «~jamais vous ne
pourrez récompenser la peine que j'ai prise pour vous autres~» &
Agree \\
Sunday & «~Je vous ai donné six jours pour travailler, je me suis
réservé le septième, et on ne veut pas me l'accorder~» & Same, with
«~et on ne veut pas me l'accorder~» & Agree \\
Blasphemy of the carters & «~Ceux qui \emph{conduisent} les
charrettes ne savent pas jurer sans \emph{y} mettre le nom de mon
fils au milieu~» & «~ceux qui \emph{mènent} les charrettes ne savent
pas jurer sans mettre le nom de mon Fils au milieu~» & Agree in
substance; two words differ \\
The two burdens & ``These are the two things which make the hand of
my Son so heavy'' & «~Ce sont les deux choses qui appesantissent
tant le bras de mon Fils~» & Agree (with the hand\slash arm
difference) \\
Spoilt harvest and potatoes & French, then repeated in
\term{patois} because Mélanie did not understand & French only, with
a one-line \term{patois} marker, «~Si la recolta se gasta\ldots~» &
\textbf{Differ in form}: the early witness prints the whole repeated
passage in \term{patois} \\
Corn, famine, children under seven, walnuts, grapes & Present, in
\term{patois} & Present, in French & Agree \\
Conversion reverses the sign & ``the stones and the rocks will
change into heaps of corn, and the potatoes will be sown of
themselves'' & «~les pierres et les rochers deviendront des monceaux
de blé et les pommes de terre seront ensemencées par les terres~» &
Agree \\
Private words to each child & Not in the recital as a stage
direction; Mélanie's recital says only that the Lady seemed to speak
to Maximin and she heard nothing, and the interrogations treat the
secret separately & Marked by parenthetical stage directions inside
the text & \textbf{Differ in form}: the received text incorporates
the secrets' moment editorially \\
Prayer & ``Do you say your prayers well?\ldots\ at least an Our
Father and a Hail Mary'' & Same, «~ne diriez-vous seulement qu'un
`Notre Père' et un `Je vous salue'~» & Agree \\
Mass, Sunday work, Lent & Present, in \term{patois} & Present, in
French & Agree \\
The Coin episode & Present, in \term{patois}, with Maximin's
recognition & Present, in French, with the same recognition & Agree
\\
Final charge & ``you will make this to be known to all my people,''
twice & «~vous le ferez passer à tout mon peuple~»;
«~faites-le bien passer à tout mon peuple~» & Agree \\
\end{stratatable}

Three conclusions follow, and only three. First, the received
custodial text is a legitimate descendant of the recital stratum:
every element in it is in the early witness, and the early witness
adds nothing the received text suppresses. Second, the received text
is a \emph{translation and harmonization}, not a transcript: the
\term{patois} has been turned into French, and the private-words
moment has been made into a stage direction. Third, at least one
substantive word differs---\term{la main} in the early witness,
\term{le bras} in the received text---and no witness available to
this project settles which the children said. Anyone who builds an
argument on the exact French of a La Salette sentence is building on
that unsettled ground.\endnote{The comparison was made line by line
between Ullathorne 1854, pp.~36--42 (French and \term{patois}
columns) and the Sanctuary's received French text quoted at \S3.3, on
2026-07-25; the results are recorded in
\work{research/source-audit.md}. The critical documentary edition of
the early La Salette records was not consulted (Appendix~A), and the
\term{main}\slash\term{bras} variant is reported as a variant
between the two reachable witnesses, not as a decision.}

\subsection{What the message says---and what it does not}

Read as a whole, the discourse is a short penitential sermon in
peasant idiom with a strikingly narrow moral field. It has exactly
two named sins: profanation of Sunday and blasphemy that puts the
name of the Son into an oath. It has one devotional instruction:
morning and evening prayer, at minimum the Our Father and the Hail
Mary. It has one complaint about practice: neglected Mass, mocking
attendance, and Lenten disregard. It has one conditional sign
structure: continued refusal brings dearth, conversion reverses it.
And it has one commission: make this pass to all her people.

What it does not contain deserves the same precision, because much
that circulates under the name of La Salette is not in it. The
public discourse contains no date, no political programme, no
dynasty, no council, no papacy, no Antichrist, no chronology of the
last things, no promise attached to any devotion, no new title, no
new doctrine, and no claim about the standing of any bishop, priest,
or religious order. The phrase for which La Salette is best known
outside the Church---the material of the ``secret''---is by
definition not part of the public discourse at all; the children
withheld it, and the mandement never judged it (\S5).

Two elements do require a word of care. The dearth prophecy is
conditional and agricultural, and it belongs to a year of real crop
failure across Europe; nothing in this study treats it as a fulfilled
or unfulfilled prediction, and the 1851 act itself does not rest on
it. The saying about children under seven is the harshest sentence
in the discourse; it is reported here because it is in the earliest
reachable recital and suppressing it would falsify the record, and
it is not offered as a prophecy to be applied to anyone, at any
time, then or now.\endnote{Editorial judgment, stated locally as
required by \work{guidance/editorial.md}. The 1846 European
potato blight and the 1846--1847 French dearth are the discourse's
immediate agricultural background; this study does not develop the
economic history and cites no source for it beyond the discourse's
own reference to the previous year's potatoes.}

\subsection{The title under which the event is received}

The Church's cult of this event grew under a title the discourse
does not use: \term{Notre-Dame Réconciliatrice}, Our Lady
Reconciler of sinners. The title is a reception, not a quotation:
the discourse speaks of a mother interceding without ceasing that
her Son not abandon a people that will not submit, and the Church
named that office \term{Réconciliatrice}. The title is attested in
Roman acts---the 1915 decree expressly preserves devotion to the
Blessed Virgin ``under the title of Reconciliatrix commonly called
\term{de la Salette}'' (\S5.7), and John XXIII in 1961 constituted
``the Blessed Virgin Mary of La Salette, invoked as Reconciliatrix
of sinners,'' heavenly patroness of an entire religious congregation
(\S9.4). Read inside Catholic doctrine, the title says exactly what
\work{Lumen gentium} 60--62 permits it to say: a created, graced,
maternal, and wholly subordinate participation in the one
reconciliation worked by Christ, who alone reconciles (2~Cor
5:18--19). It does not make Mary a second source of reconciliation,
and no act cited in this study says or implies that it
does.\endnote{Doctrinal frame: \work{Lumen gentium} 60--62 and
\work{Catechism of the Catholic Church} 970 (Mary's saving influence
``flows forth from the superabundance of the merits of Christ,
rests on his mediation, depends entirely on it, and draws all its
power from it''). The Roman texts using the title are quoted at
their own loci in \S5.7 and \S9.4.}
